% TOP_PROBLEM: {"id": 30006970, "title": "Ruzsa-Szemeredi (6,3) Extremal-Order Problem", "category_id": 2, "category": {"id": 2, "name": "combinatorics", "display_name": "Combinatorics", "description": "Counting problems, graph theory, discrete structures.", "slug": "combinatorics", "order_index": 2, "created_at": "2026-07-31T15:26:25.671Z"}, "status": "open"}
% FIRST_POSTED: 2026-09-27
% !TeX program = lualatex
% ATTEMPT_STATUS: unresolved
\documentclass[11pt]{article}
\usepackage[margin=25mm]{geometry}
\usepackage{fontspec}
\setmainfont{Latin Modern Roman}
\usepackage{amsmath,amssymb,mathtools}
\usepackage{unicode-math}
\setmathfont{Latin Modern Math}
\usepackage{xurl,hyperref}
\hypersetup{colorlinks=true,urlcolor=blue}
\setcounter{secnumdepth}{0}
\setlength{\parindent}{0pt}
\setlength{\parskip}{5pt}
\setlength{\emergencystretch}{3em}
\begin{document}
\section*{\texorpdfstring{Ruzsa-Szemeredi (6,3) Extremal-Order Problem}{Ruzsa-Szemeredi (6,3) Extremal-Order Problem}}\label{30006970}
\subsection{Definitions and mathematical statement}
A simple three-uniform hypergraph on an $n$-element vertex set has distinct three-element subsets as edges. Define
\[
 f(n,6,3)=\max\{|E(H)|:\ |V(H)|=n,
 \ \forall U\subseteq V(H),\ |U|\le6,
 \ \#\{e\in E(H):e\subseteq U\}\le2\}.
\]
Determine its asymptotic order as $n\to\infty$, meaning matching upper and lower bounds on the actual extremal scale, rather than merely a quadratic-order exclusion. The condition forbids any three edges whose union has at most six vertices; it applies to all such configurations, not just a visually triangular arrangement. No multiplicities are allowed among edges.
\par\smallskip\textit{Formulation and concept references:} \href{https://arxiv.org/html/2003.02754}{[S1]}.

\subsection{Short English statement}
How many triples can a large vertex set support when no collection of at most six vertices contains three of them?
\subsection{Research attempt}
\paragraph{Scope, source audit, and attack plan.}
This attempt by GPT-6 Astra Ultra was prepared on 27 September 2026. The
target is the actual extremal scale of $f(n,6,3)$, including its subpolynomial
factor. Establishing only $n^{2-o(1)}\leq f(n,6,3)=o(n^2)$ does not answer
it. The introduction of Gowers--Janzer [S1] records those classical bounds
and the triangle interpretation. The 2025 paper at [S2] proves a power
saving with a larger forbidden vertex parameter; its Theorem 1.3 concerns
$f(n,e+\lfloor\log_2 e\rfloor+38,e)$. Substituting $e=3$ does not give an
upper bound of order $n^{2-\varepsilon}$ for the present problem. Such a
bound would contradict the construction below. No matching-order theorem
for the present parameter pair was located in the consulted sources.

The proof track first removes nonlinearity exactly, then examines a spectral
obstruction to pseudorandomness in the associated graph. The construction
track tests an arithmetic family against \emph{all} forbidden three-edge
patterns. The familiar progression-free construction is included with a
complete verification to calibrate possible upper bounds; it is not a claim
of a new extremal construction. The remaining issue is quantitative global
structure, rather than detecting one local triangular pattern.

\paragraph{Lemma 494.1: an exact decomposition, including the non-linear edges.}
Let $L(n)$ be the maximum number of edges in a linear $(6,3)$-free
three-uniform hypergraph on $n$ vertices; linear means that different edges
intersect in at most one vertex. Put $L(0)=L(1)=L(2)=0$. Then
\begin{equation}\label{a494:decomposition}
 f(n,6,3)=\max_{0\leq k\leq\lfloor n/4\rfloor}
                  \bigl(2k+L(n-4k)\bigr).
\end{equation}
In particular $L(n)\leq f(n,6,3)\leq L(n)+n/2$.

\emph{Proof.} Suppose two distinct edges $e_1,e_2$ share two vertices.
Their union $W$ has four vertices. Any third edge meeting $W$ would have
union of size at most $4+3-1=6$ with the first two edges, which is forbidden.
Thus these two edges form an isolated component on exactly four vertices.
Every such component is disjoint from every other one. After removing its
$k$ four-vertex components of this kind, the hypergraph has $n-4k$ vertices
and is linear. This proves the upper bound in \eqref{a494:decomposition}.

Conversely, take a linear admissible hypergraph on $n-4k$ vertices and
adjoin $k$ disjoint copies of the two triples $\{1,2,3\}$ and
$\{1,2,4\}$. Three edges in the old component remain admissible. Three
edges meeting more than one component either include two edges whose union
has at least four vertices and a disjoint triple, or are three disjoint
triples. Their union therefore has at least seven vertices. A new
four-vertex component has only two edges. The union is admissible, proving
the other inequality. Finally $L(n-4k)\leq L(n)$ by adjoining isolated
vertices, and $2k\leq n/2$.\hfill$\square$

This reduction avoids a tempting but false preliminary assertion that every
admissible hypergraph is linear. The two triples just displayed are an
explicit counterexample to that assertion. Linearity is obtained after a
controlled loss, and the exact formula records that loss.

\paragraph{Lemma 494.2: the triangle model is equivalent after that reduction.}
The quantity $L(n)$ equals the maximum number of triangles in an $n$-vertex
simple graph in which every edge belongs to exactly one triangle. Isolated
vertices are allowed.

\emph{Proof.} Replace each hyperedge of a linear admissible hypergraph
by its three pairs, producing a graph $G$. Each pair occurs in exactly one
of the designated triangles. Suppose $abc$ is another triangle. The
hyperedges containing $ab,bc,ca$ must be distinct: if one contained two of
these pairs, it would be the triple $abc$, and uniqueness of pairs would
force all three designated hyperedges to be this same triple. The three
distinct hyperedges have union inside $\{a,b,c\}$ plus their three remaining
vertices, a forbidden set of size at most six. There are no extra triangles.

For the converse, take the triangles of such a graph as hyperedges. They
are linear because an edge of the graph cannot lie in two triangles.
For three distinct linear triples, inclusion--exclusion gives
\[
 |e_1\cup e_2\cup e_3|
 =9-\sum_{i<j}|e_i\cap e_j|+|e_1\cap e_2\cap e_3|.
\]
The union has at most six vertices only if all three pairwise intersections
are singletons and the triple intersection is empty. The three intersection
vertices are then distinct and form an additional graph triangle. Each of
its edges is already in one of the three original triangles, a contradiction.
This proves both directions.\hfill$\square$

If the graph has $m$ triangles, it has exactly $3m$ edges. Therefore
$m\leq\binom n2/3$. Every vertex neighborhood induces a perfect matching:
each neighbor $u$ of $v$ is paired with the unique third vertex of the
triangle containing $uv$, and any further adjacency within the neighborhood
would put $uv$ in another triangle. In particular all degrees are even.
These local facts will be used as checks, not as a substitute for the
global extremal estimate.

\paragraph{Construction track: a fully checked arithmetic family.}
Let $A\subseteq\{1,\ldots,N\}$ have no nonconstant three-term arithmetic
progression. Use three disjoint labelled vertex classes
\[
 X=\{1,\ldots,N\},\quad Y=\{1,\ldots,2N\},\quad
 Z=\{1,\ldots,3N\},
\]
and for every $(x,a)\in X\times A$ insert the triple
\begin{equation}\label{a494:arithmetic}
                   (x,x+a,x+2a)\in X\times Y\times Z.
\end{equation}
The labels distinguish the classes even when integer values agree. There
are exactly $N|A|$ distinct triples on $6N$ vertices. Any two coordinates
determine $x$ and $a$, so the hypergraph is linear.

To check the forbidden configuration condition without assuming that it is
already a triangular one, consider the shadow graph from Lemma 494.2.
Every triangle must use one vertex from each class. Its three pair edges
imply that
\[
 a=y-x\in A,\qquad b=z-y\in A,\qquad c=(z-x)/2\in A.
\]
Hence $a+b=2c$. Progression-freeness implies $a=b=c$, and the triangle is
one of \eqref{a494:arithmetic}. Each pair is in just one such triangle.
Lemma 494.2 now checks every collection of three hyperedges, including
degenerate intersection patterns. We have proved the exact finite bound
\begin{equation}\label{a494:r3}
 f(6N,6,3)\geq L(6N)\geq N r_3(N),
\end{equation}
where $r_3(N)$ is the largest size of a progression-free subset of $[N]$.
This is a lower bound; the argument supplies no reverse inequality comparing
arbitrary admissible hypergraphs with subsets of an interval.

Here is an explicit version of the geometric digit construction behind a
large $A$. For integers $k,d\geq2$, partition $\{0,\ldots,k-1\}^d$ by
the value of $\sum_i v_i^2$. One fiber $V$ has
\[
 |V|\geq\frac{k^d}{d(k-1)^2+1}.
\]
On a fixed sphere, $u+w=2v$ implies $u=v=w$: expanding the squared norms
gives $\|u-w\|^2=0$. Map the fiber into integers by
\[
                 v\longmapsto 1+\sum_{i=0}^{d-1}v_i(2k)^i.
\]
Its values lie in $[1,(2k)^d]$. In an equation between the sum of two
such values and twice a third, all digit sums are at most $2k-2<2k$.
There are no carries, so equality of integers gives equality of digit
vectors. The image is therefore progression-free. In particular, writing
$N_0=(2k)^d$,
\begin{equation}\label{a494:digitbound}
        \frac{|A|}{N_0}\geq
        \frac{2^{-d}}{d(k-1)^2+1}.
\end{equation}

For arbitrary sufficiently large $N$, put $d=\lfloor\sqrt{\log N}\rfloor$
and $k=\lfloor N^{1/d}/2\rfloor$. Then $(2k)^d\leq N$ and, for large
$N$, $k\geq N^{1/d}/3$. Consequently
\[
 |A|\geq\frac{N3^{-d}}{dN^{2/d}+1}
                 \geq N\exp(-C\sqrt{\log N})
\]
for an absolute $C$. All logarithms here are natural. Using
$N=\lfloor n/6\rfloor$ and adding isolated vertices yields
\begin{equation}\label{a494:lower}
                  f(n,6,3)\geq
                  n^2\exp(-C'\sqrt{\log n})
\end{equation}
for all sufficiently large $n$. The exponent loss is $o(\log n)$, so any
proposed upper bound $O(n^{2-\varepsilon})$ with fixed $\varepsilon>0$
fails this stress test. Constants may change between displays.

\paragraph{Proof track: why qualitative removal does not determine the scale.}
The triangle removal lemma asserts that for every $\eta>0$ there is
$\delta(\eta)>0$ such that a graph with at most $\delta(\eta)n^3$
triangles can be made triangle-free by deleting at most $\eta n^2$ edges.
This is the external theorem used here; [E1] is a primary proof. In the
graph of Lemma 494.2, exactly $m$ deletions are necessary and sufficient:
the $m$ triangles are edge-disjoint, and deleting one edge from each kills
all of them.

Fix $\varepsilon>0$. If $m\geq\varepsilon n^2$, use removal with
$\eta=\varepsilon/2$. Since $m\leq n^2/6$, we have
$m/n^3\leq1/(6n)<\delta(\varepsilon/2)$ for sufficiently large $n$.
Removal would require at most $\varepsilon n^2/2<m$ deletions, a
contradiction. Thus $L(n)=o(n^2)$, and Lemma 494.1 gives the same for
$f(n,6,3)$. This recovers a classical conclusion, not the requested
matching asymptotic estimate.

The quantitative dependence has not vanished. If $\alpha=m/n^2>0$, the
same reasoning only forces
\[
                      \delta(\alpha/2)<\alpha/n
\]
with the convention using a non-strict triangle threshold. To turn this
into a proposed bound on $\alpha$, one must supply and invert an actual
function $\delta$. Replacing it by a polynomial without proof would be
especially damaging: this would conflict with the family
\eqref{a494:lower}. The procedure cannot simply infer the desired rate
from the fact that a positive $\delta$ exists.

\paragraph{Lemma 494.3: dense triangle systems force a large secondary eigenvalue.}
Let $G$ have $n\geq2$ vertices, every edge in exactly one triangle, and
$m>0$ triangles. Let $\lambda_1\geq\cdots\geq\lambda_n$ be the real
eigenvalues of its adjacency matrix, and set
$\Lambda=\max_{2\leq i\leq n}|\lambda_i|$. Then
\begin{equation}\label{a494:spectral}
                  \Lambda\geq\frac{36m^2}{n^3}-1.
\end{equation}

\emph{Proof.} An adjacency matrix has
$\operatorname{tr}(A^2)=2|E(G)|=6m$ and
$\operatorname{tr}(A^3)=6\#K_3(G)=6m$, by counting closed walks of
length two and three. Its largest eigenvalue is at least the average
degree, by the Rayleigh quotient at the all-ones vector:
$\lambda_1\geq6m/n$. For $i\geq2$,
$\lambda_i^3\geq-\Lambda\lambda_i^2$. Consequently
\[
 6m=\sum_i\lambda_i^3
      \geq\lambda_1^3-\Lambda\sum_{i\geq2}\lambda_i^2
      \geq\lambda_1^3-6m\Lambda.
\]
Divide by $6m$ and use the Rayleigh bound.\hfill$\square$

For $m=\alpha n^2$, this gives $\Lambda\geq36\alpha^2 n-1$.
In particular a hypothetical dense family with every nonprincipal
eigenvalue $o(n)$ is impossible. This is a structural restriction with an
elementary proof, but it is consistent with the arithmetic construction:
the latter is highly structured and its density tends to zero. The estimate
does not yield a contradiction for a slowly decreasing $\alpha(n)$.

A contemplated next step was to use a large eigenvector to find a denser
induced subsystem and iterate. Here is the first missing assertion: a
spectral cut must retain enough \emph{whole triangles}, with a quantitatively
controlled density gain, to repeat the argument. A cut can retain edges
whose unique completing vertex was discarded. Those edges must then be
removed; their number is not controlled by \eqref{a494:spectral}. A density
increment for graph edges by itself is insufficient. No such uniform
triangle-preserving increment is proved in this attempt.

\paragraph{Finite verification and adversarial checks.}
The companion Python script exhaustively checks the inclusion--exclusion
criterion on 56,155 linear three-edge configurations using vertex sets of
size at most nine. It also tests all 260 progression-free subsets of
$[N]$ for $1\leq N\leq8$, constructs their hypergraphs, checks pair
uniqueness, and explicitly checks the union of every three edges. One
test is $N=8$, $A=\{1,2,4,5\}$, giving 32 triples on 48 available
vertices. Four small sphere constructions are checked separately for
progression-freeness. These computations test implementations of the
proved lemmas; they do not approximate the unknown extremal order.

Three boundary errors were specifically excluded. The union condition is
``at most six'', so non-linear degeneracies cannot be omitted. Empty vertex
classes or extra isolated vertices do not change admissibility. Finally,
integer division by two in the shadow construction is legitimate because
an $XZ$ edge supplies an actual integer step in $A$; replacing the interval
by an even cyclic group would invalidate the pair-uniqueness argument.

\paragraph{Outcome and next concrete targets.}
\textbf{UNRESOLVED.} The independently checked work gives the exact
non-linear decomposition, the graph equivalence, a fully verified
progression-free family, and the spectral inequality
\eqref{a494:spectral}. It reproduces the known broad bounds without
matching their subpolynomial factors. No new extremal-order theorem is
claimed. The first gap in the pursued upper-bound route is the
triangle-preserving density-increment statement described above.

Further work could quantify the loss of triangles under spectral cuts,
test such a bound against the digit constructions, and seek a decomposition
whose parts preserve unique triangle completion. The construction route
would instead need larger progression-free or non-arithmetic systems on
the same vertex budget. Even matching \eqref{a494:lower} on a logarithmic
exponent scale would require a proof; it would not automatically identify
the extremal function up to constant factors.

\subsection{Sources}
\begin{itemize}
\item[S1]Gowers and Janzer, Generalizations of the Ruzsa-Szemeredi and rainbow Turan problems for cliques (2020).
\url{https://arxiv.org/html/2003.02754}.
\textit{Formulation source.}
\item[S2]Status source for Ruzsa-Szemeredi (6,3) Extremal-Order Problem.
\url{https://people.math.ethz.ch/~sudakovb/Brown-Erdos-Sos-power-saving.pdf}.
\textit{Further reference; not a proof endorsement.}
\item[E1]Jacob Fox, A new proof of the graph removal lemma,
Annals of Mathematics 174 (2011), 561--579.
\url{https://annals.math.princeton.edu/2011/174-1/p17}.
\url{https://arxiv.org/abs/1006.1300}.
\textit{Primary source for the removal theorem; consulted 27 September 2026.}

\end{itemize}
\end{document}
